\documentclass{article}
\usepackage{graphicx}
\usepackage{amsfonts}
\usepackage{amsmath}
\usepackage{amsthm}
\usepackage{tikz}

\theoremstyle{definition}
\newtheorem{theorem}{Theorem}[section]
\newtheorem{corollary}[theorem]{Corollary}
\newtheorem{proposition}[theorem]{Proposition}
\newtheorem{lemma}[theorem]{Lemma}
\newtheorem{definition}[theorem]{Definition}

\theoremstyle{remark}
\newtheorem{remark}[theorem]{Remark}
\newtheorem{example}[theorem]{Example}
\newtheorem{claim}[theorem]{Claim}

\title{Sequences}
\author{Sarah Liu}
\date{January 2026}

\begin{document}

\maketitle

\section{Monotonic and Bounded Sequence}

\begin{definition}[Monotonic]
    A sequence is \textit{increasing} when $\forall n\in\mathbb{N}, a_n<a_{n+1}$.
    Or equivalently, $\forall n,m\in\mathbb{N}, n<m\implies a_n<a_m$.

    A sequence is \textit{decreasing} when $\forall n\in\mathbb{N}, a_n>a_{n+1}$.
\end{definition}

\begin{definition}[Bounded]
    Let $\{a_n\}_{n=0}^{\infty}$ be a sequence.
    The sequence is bounded if there exists $A,B\in\mathbb{R}$ such that
    \[
    \forall n\in\mathbb{N},\quad A\leq a_n\leq B
    \]
\end{definition}

\begin{definition}[Convergent]
    Let $\{a_n\}_{n=0}^{\infty}$ be a sequence.
    The sequence is convergent if there exists $L\in\mathbb{R}$ such that
    \[
    \forall\varepsilon >0, \;\exists n\in\mathbb{N}\; \text{s.t.}\;\forall n\in\mathbb{N},\; n\geq n_0\implies|a_n-L|<\varepsilon
    \]
\end{definition}

\begin{definition}
    Let $\{a_n\}_{n=0}^{\infty}$ be a sequence.
    The sequence is divergent to $\infty$ if 
    \[
    \forall M >0, \;\exists n\in\mathbb{N}\; \text{s.t.}\;\forall n\in\mathbb{N},\; n\geq n_0\implies|a_n-L|<\varepsilon
    \]
\end{definition}

\begin{theorem}
    Let $\{a_n\}_{n=0}^{\infty}$ be a sequence.
    If the sequence is convergence, then it is bounded.
\end{theorem}

\begin{theorem}[the Monotone Convergence Theorem for Sequence]
    If a sequence is eventually monotonic and bounded, then it is convergent.
\end{theorem}

\section{the Big Theorem}

\begin{theorem}
    \[
    \ln n << n^a << c^n << n! <<n^n
    \]
    where $a>0$, $c>1$.
\end{theorem}

\begin{lemma}
    \[
    \lim_{n\rightarrow\infty} \frac{\ln n}{n^a}=0
    \]
\end{lemma}

\begin{proof}
    \[
    \lim_{n\rightarrow\infty} \frac{\ln n}{n^a}= \lim_{n\rightarrow\infty} \frac{\frac{1}{n}}{an^{a-1}}= \lim_{n\rightarrow\infty} \frac{1}{an^a}=0
    \]
\end{proof}

\begin{lemma}
    \[
    \lim_{n\rightarrow\infty} \frac{n^a}{c^n}=0
    \]
\end{lemma}

\begin{proof}
    \[
    \lim_{n\rightarrow\infty} \frac{n^a}{c^n}= \lim_{n\rightarrow\infty}\frac{an^{a-1}}{\ln c\cdot c^n}=\dots =\lim_{n\rightarrow\infty}\frac{a!}{(\ln c)^a\cdot c^n}=0
    \]
\end{proof}

\begin{lemma}
    \[
    \lim_{n\rightarrow\infty} \frac{c^n}{n!}=0
    \]
\end{lemma}

\begin{proof}
    For this proof, we will use induction on $n$.
    Assume the conclusion holds for $n$, then 
    \[
    \lim_{n\rightarrow\infty} \frac{c^{n+1}}{{n+1}!} = \lim_{n\rightarrow\infty} \frac{c}{{n+1}}\frac{c^{n}}{{n}!} = \lim_{n\rightarrow\infty} \frac{c}{{n+1}}\cdot\lim_{n\rightarrow\infty} \frac{c^{n}}{{n}!}=0\cdot 0=0
    \]
\end{proof}

\begin{proof}[Alternative]
    Fix $c>1$.
    Call $p_n=\frac{c^n}{n!}$.
    Because $p_{n+1}=\frac{c}{n+1}p_n$, then when $n\geq c$, we have $p_{n+1}<p_n$. So the sequnce is eventually decreasing.
    Also, $\{p_n\}$ is bounded below by 0.
    By MCT, we know that $\{p_n\}$ is convergent.
    \[
    \lim_{n\rightarrow\infty} p_{n+1}=\lim_{n\rightarrow\infty} \frac{c}{{n+1}}\cdot\lim_{n\rightarrow\infty} p_n
    \]
    Let $L=\lim_{n\rightarrow\infty} p_{n+1}$, then $L\cdot0=L$. So $L=0$.
\end{proof}

\begin{lemma}
    \[
    \lim_{n\rightarrow\infty} \frac{n!}{n^n}=0
    \]
\end{lemma}

\begin{proof}
    \[
    0<\frac{n!}{n^n}=\frac{1}{n}\cdot \frac{2}{n}\cdot \frac{3}{n}\dots \frac{n-1}{n}\cdot \frac{n}{n}=\frac{1}{n}\cdot \left(\frac{2}{n}\cdot \frac{3}{n}\dots \frac{n-1}{n}\cdot \frac{n}{n}\right)\leq \frac{1}{n}
    \]
    Then because $\lim_{n\rightarrow\infty}0=\lim_{n\rightarrow\infty}1/n=0$, with the squeeze theorem, we know that 
    \[
    \lim_{n\rightarrow\infty} \frac{n!}{n^n}=0
    \]
\end{proof}

\end{document}